\section{The complete relative endomorphism spectrum}
\label{sec:relative-spectrum}

The complementary correction in Section~\ref{sec:projected-he}
provides actual connections for every small matrix parameter,
even before solving the remaining finite equation.
The metric isometry $U_p(z)=e^{s_p(z)/2}$ transports the
Chern connection of $(\dbar_0+b_p(z),h_0e^{s_p(z)})$ to a
unitary connection $\mathcal A_p(z)$ for the fixed metric
$h_0$. Its moment is exactly
\[
                 \mathrm i\Lambda F_{\mathcal A_p(z)}=\Psi_p(z).
\]
On the real Hilbert space of trace-free Hermitian endomorphisms
let
\begin{equation}\label{eq:actual-rough-form}
 \mathcal L_{p,z}=D_{\mathcal A_p(z)}^*D_{\mathcal A_p(z)},
 \qquad q_{p,z}(u)=\|D_{\mathcal A_p(z)}u\|_2^2.
\end{equation}
The form domain is $H^1$ on each compact smooth fibre;
the operator has its usual elliptic domain. On the parallel
space $\mathfrak h$ define
\begin{equation}\label{eq:finite-bundle-energy}
 \begin{split}
 E_{p,z}(S)={}&2c_+(p)\|S_+A-AS_-\|_{\rm HS}^2\\
             &+2c_-(p)\|S_-B-BS_+\|_{\rm HS}^2.
 \end{split}
\end{equation}
Its mass is $M$ in \eqref{eq:split-mass}. This form uses
all Hermitian matrix directions, including noncommuting
ones. Denote its generalized eigenvalues, including zeros,
by $\mu_1\le\cdots\le\mu_7$.

\begin{theorem}\label{thm:complete-bundle-spectrum}
Uniformly for $p\in\mathcal B^\circ$ near zero and $z$
in a fixed small matrix ball, the whole form satisfies
\begin{equation}\label{eq:whole-bundle-form}
 \begin{split}
 (1-C|z|)\{E_{p,z}(S)+q_{p,0}(w)\}
 &\le q_{p,z}(S+w)\\
 &\le(1+C|z|)\{E_{p,z}(S)+q_{p,0}(w)\},
 \end{split}
\end{equation}
for $S\in\mathfrak h$, $w\perp\mathfrak h$ in $H^1$.
The exact kernel consists of the constant sections
\begin{equation}\label{eq:actual-stabilizer-kernel}
 \{S\in\mathfrak h:S_+A=AS_-,\quad S_-B=BS_+\}.
\end{equation}
The ordered eigenvalues of $\mathcal L_{p,z}$, including
zeros, satisfy
\begin{equation}\label{eq:seven-eigenvalues}
 (1-C|z|)\mu_j\le\lambda_j\le(1+C|z|)\mu_j
       \quad(1\le j\le7),\qquad\lambda_8\ge c>0.
\end{equation}
The split projection $\Pi_p$ identifies the whole
seven-dimensional spectral space with $\mathfrak h$,
and its inverse differs from the constant inclusion by
$O(|z|)$ in $L^2$ operator norm. All these conclusions
allow arbitrary ratios among the positive $\mu_j$ and
changing kernel dimension.
At a zero of $\Psi_p$, this is the complete small spectrum
of the linearization of
$h^{-1}\partial_th=-2\mathrm i\Lambda F_h$.
\end{theorem}

\begin{proof}
The transported $(0,1)$ coefficient is exactly
\[
 U_p b_p U_p^{-1}-(\dbar_0U_p)U_p^{-1}.
\]
The graph and parameter bounds for $s_p$, together with
\eqref{eq:controlled-harmonic-gauge}, give
\begin{equation}\label{eq:connection-expansion}
 \begin{gathered}
 \mathcal A_p(z)=\mathcal A_{0,p}+a_{1,p}(z)+\mathcal N_p(z),
 \qquad a_{1,p}(z)=\eta_p(z)-\eta_p(z)^*,\\
 \|a_p(z)\|_\infty\le C|z|,
 \quad\|\mathcal N_p(z)\|_\infty\le C|z|^2,
 \quad\|d_z\mathcal N_p(z)\|_\infty\le C|z|,
 \end{gathered}
\end{equation}
where $a_p=\mathcal A_p-\mathcal A_{0,p}$.
These maps are $K$-equivariant. For
\[
 \delta_zS=(S_+A-AS_-,\ S_-B-BS_+)
\]
the infinitesimal action of $\mathrm iS$ is
$\mathrm i\delta_zS$. Differentiating equivariance gives
\[
 [\mathrm iS,\mathcal N_p(z)]
             =d_z\mathcal N_p(z)[\mathrm i\delta_zS].
\]
The positive lower bounds for $c_\pm$ and fixed volume
therefore imply the crucial directional estimate
\begin{equation}\label{eq:relative-commutator}
 \|[\mathcal N_p(z),S]\|_2
                         \le C|z|\,E_{p,z}(S)^{1/2}.
\end{equation}
It holds also when $E_{p,z}(S)=0$, in which case the
commutator vanishes exactly. This estimate differentiates
the compact action, so it does not require a holomorphic
connection remainder.

The type components and the two off-diagonal blocks are
orthogonal. Hence
\[
 \|[a_{1,p},S]\|_2^2
          =2\|[\eta_p,S]\|_2^2=E_{p,z}(S).
\]
Moreover $a_{1,p}$ is coclosed for $D_0$: the $(0,1)$
component is $\dbar_0$-coclosed by harmonicity, the other
component is its adjoint, and the remaining adjoints have
the wrong bidegrees. Since $S$ is parallel,
\begin{equation}\label{eq:linear-orthogonality}
       \langle D_0w,[a_{1,p},S]\rangle_2=0.
\end{equation}
The split section gap and \eqref{eq:connection-expansion}
also give
\[
 \|[a_p,w]\|_2\le C|z|\|w\|_2
                         \le C'|z|q_{p,0}(w)^{1/2}.
\]
Now expand
\[
 D_{\mathcal A_p(z)}(S+w)
     =D_0w+[a_{1,p},S]+[\mathcal N_p,S]+[a_p,w].
\]
By \eqref{eq:linear-orthogonality} the squared norm of
the first two terms is exactly $q_{p,0}(w)+E_{p,z}(S)$.
The other two have norm at most $C|z|$ times its square
root. Squaring proves \eqref{eq:whole-bundle-form}, first
on smooth sections and then on the entire form domain.

The lower bound forces $w=0$ and $E_{p,z}(S)=0$ for
any vector in the kernel. Conversely these conditions
annihilate both commutators, including the nonlinear
one by \eqref{eq:relative-commutator}. This proves
\eqref{eq:actual-stabilizer-kernel}.
The comparison form $E_{p,z}(S)+q_{p,0}(w)$ is an
orthogonal direct sum. Its first seven eigenvalues are
$\mu_1,\ldots,\mu_7$, because $\mu_7\le C|z|^2$ while
the complementary gap is fixed. Min--max applied to
\eqref{eq:whole-bundle-form} proves
\eqref{eq:seven-eigenvalues}, with the exact zero
multiplicities already established.

On the entire small spectral space the energy is at most
$C|z|^2$ times the norm. If $v=S+w$ belongs to it, the
lower comparison and split gap yield
$\|w\|_2\le C|z|\|v\|_2$. Thus $\Pi_p$ is injective
there and hence an isomorphism onto the seven-dimensional
space $\mathfrak h$, with the asserted inverse estimate.
This is a whole-space identification, without a claim
about individual nearly coincident eigenlines.

At $\Psi_p(z)=0$ the corrected connection is HE.
The section Kähler identity gives
$\mathcal L_{p,z}=2\dbar_{\mathcal A_p(z)}^*
\dbar_{\mathcal A_p(z)}$. The infinitesimal curvature
moment for the variation $h\mapsto he^v$ is
$\dbar^*\dbar v$. The prescribed metric flow therefore
linearizes to $\partial_tv=-\mathcal L_{p,z}v$.
Away from that zero set the theorem concerns the rough
operator \eqref{eq:actual-rough-form}; it does not
identify a Hessian at a noncritical metric.
\end{proof}

A useful second reconstruction minimizes energy at a
fixed split projection. For each $S\in\mathfrak h$,
let $w(S)\perp\mathfrak h$ minimize $q_{p,z}(S+w)$.
Complementary coercivity gives a unique linear minimizer.
The expansion in the proof shows
\[
 |q_{p,z}(S,w)|\le C|z|E_{p,z}(S)^{1/2}q_{p,0}(w)^{1/2}.
\]
Testing its equation with $w(S)$ gives
\begin{equation}\label{eq:variational-reconstruction}
 q_{p,0}(w(S))^{1/2}\le C|z|E_{p,z}(S)^{1/2}.
\end{equation}
Its mass is $M(S,S)+\|w(S)\|_2^2$, and the minimized
energy is relatively equivalent to $E_{p,z}$.
This variational reconstruction vanishes on the exact
kernel. It need not be the inverse of the spectral
projection used in Theorem~\ref{thm:complete-bundle-spectrum}.

\subsection{The quadratic term of the actual finite moment}

\begin{proposition}\label{prop:quadratic-moment}
The map in \eqref{eq:finite-actual-moment} satisfies
\begin{equation}\label{eq:moment-taylor}
 \Psi_p(z)=\mu_{2,p}(z)+O(|z|^3),\qquad
 d_z\Psi_p(z)=d_z\mu_{2,p}(z)+O(|z|^2),
\end{equation}
where
\begin{equation}\label{eq:quadratic-moment}
 \mu_{2,p}(A,B)=\frac1V
 \bigl(c_+AA^*-c_-B^*B,\ c_-BB^*-c_+A^*A\bigr).
\end{equation}
The estimates are uniform in the full period domain.
\end{proposition}

\begin{proof}
Pair the contracted differentiated curvature terms of
$M_b$ with a parallel block matrix. The Kähler adjoint
identities and integration by parts show that their
projection is zero. Consequently
\[
 \Pi_p M_b=\Pi_p\mathrm i\Lambda(-b\wedge b^*-b^*\wedge b).
\]
Replacing $b$ by $\eta_p$ changes this by $O(|z|^3)$.
In a unitary coframe with
$\omega=\mathrm i\sum_j\theta^j\wedge\bar\theta^j$,
if $\beta=\sum_jC_j\bar\theta^j$, then
\[
 \mathrm i\Lambda(-\beta\wedge\beta^*-\beta^*\wedge\beta)
                                      =\sum_j[C_j,C_j^*].
\]
The two blocks of \eqref{eq:harmonic-arrow-form} therefore
give the diagonal expression
\[
 \begin{pmatrix}
 |\eta_p^+|^2AA^*-|\eta_p^-|^2B^*B&0\\
 0&|\eta_p^-|^2BB^*-|\eta_p^+|^2A^*A
 \end{pmatrix}.
\]
Integration supplies the coefficients $c_\pm(p)$ and
projection divides by $V$. The total trace is zero,
so no further scalar correction occurs.
The explicit expansion of
\eqref{eq:finite-chern-change} also gives
\[
 \|\mathcal F_p(s,z)-M_b-P_ps\|_\infty
       \le C(|z|+\|s\|_{\mathcal G})\|s\|_{\mathcal G}.
\]
At $s=s_p(z)=O(|z|^2)$ this is $O(|z|^3)$, and
$\Pi_pP_ps=0$. This proves the first estimate.
The uniform differentiated Taylor expansions give the
second estimate.
\end{proof}

An additive cubic error in \eqref{eq:moment-taylor}
can dominate an arbitrarily small quadratic component.
We therefore solve the exact equation in the next
section, instead of inferring its zeros from this
expansion. For the spectral form itself, replacing
$c_\pm(p)$ by $c_{\pm,0}$ is already legitimate:
it changes \eqref{eq:finite-bundle-energy} by a relative
factor $1+O(e(p))$, where
\[
 e(p)=\max_\pm|c_\pm(p)/c_{\pm,0}-1|\longrightarrow0.
\]
Thus the relative spectral error is $O(|z|)+O(e(p))$
with these fixed central coefficients, for every joint
base and bundle path in the stated domain.
